%! @@TITLE@@ — Unit @@UNITINT@@, Lesson @@LESSONID@@ — Slides (Beamer)
% SLIDES-FIRST shape (course shape 2026-09-25): the deck IS the lesson. Its printed handout
% (lessonYY_slides.pdf, 3-up with a notes column) is the student's notes — definitions printed in
% full, examples worked, and every "Now you try" left open with the notes column as work space.
%   title → targets → warm-up → { definition → example → now you try } × 3–4 → wrap-up → homework
% Answers are revealed on a click with \reveal{...} (algebra2-beamer): the projector shows them,
% the printed handout (a Beamer handout-mode compile) never does. \pause inside a worked example
% is fine — the handout collapses it and prints the full worked example.
\documentclass[aspectratio=169, 11pt]{beamer}
\usepackage{@@PREFIX@@-beamer}   % \forestheader{}, \sectionlabel[color]{}, \reveal[n]{}

\begin{document}

% TITLE SLIDE (CourseName is not defined in beamer, so the name is literal)
{
\setbeamercolor{background canvas}{bg=forest}
\begin{frame}[plain]
  \vspace{2.8cm}\hspace{0.55cm}%
  \begin{minipage}{0.9\textwidth}
    {\color{goldacc}\bfseries\small LESSON @@LESSONID@@}\\[0.5cm]
    {\color{white}\bfseries\LARGE @@TITLE@@}\\[1.2cm]
    {\color{forestmid}\small @@COURSENAME@@~$\cdot$~Unit @@UNITINT@@}
  \end{minipage}
\end{frame}
}

% LEARNING TARGETS
\begin{frame}[t, plain]
  \forestheader{Today's targets}
  \sectionlabel{What you will be able to do}
  \begin{itemize}\setlength{\itemsep}{6pt}
    \item TODO: target, naming the vocabulary in \textbf{bold}.
  \end{itemize}
  \vspace{4pt}
  \begin{block}{How today runs}
    Warm-up (5) $\to$ definitions, examples, and \textbf{now you try} (40) $\to$ wrap-up (5)
    $\to$ \textbf{start the homework} (10).
  \end{block}
\end{frame}

% WARM-UP — two or three spiral-review items; answers on a click. The last one leaves the question
% the first definition answers.
\begin{frame}[t, plain]
  \forestheader{Warm-Up}
  \sectionlabel{5 minutes --- in the notes column, alone}
  \begin{enumerate}\setlength{\itemsep}{7pt}
    \item TODO item. \reveal{\; \textbf{TODO answer}}
  \end{enumerate}
\end{frame}

% ── ONE CYCLE PER IDEA — repeat these three frames 3–4 times ──────────────────
% DEFINITION — complete printed sentences, term in bold. Nothing to fill in: students annotate in
% the notes column. A pre-drawn display where one helps.
\begin{frame}[t, plain]
  \forestheader{Definition: TODO term}
  \sectionlabel{TODO idea N}
  \begin{block}{TODO term}
    TODO: the definition or general form, as a complete sentence.
  \end{block}
  % TODO: a pre-drawn display (graph, table), or the two things students conflate side by side.
\end{frame}

% EXAMPLE — the "I do": worked in full, every step with its reason. \pause between steps if wanted.
\begin{frame}[t, plain]
  \forestheader{Example N: TODO}
  \sectionlabel{Watch --- I do this one}
  TODO: the problem.
  % TODO: the worked solution, step by step.
\end{frame}

% NOW YOU TRY — the "you do": one or two problems of the SAME shape as the example, fresh numbers.
% Answer (and the key step) inside \reveal{}. The last cycle's now-you-try is the CRUX — the case
% where the two answers disagree: flag it \sectionlabel[redacc]{...}.
\begin{frame}[t, plain]
  \forestheader{Now you try}
  \sectionlabel[goldacc]{In the notes column --- then check}
  TODO: the problem(s).
  \vspace{6pt}
  \reveal{%
    \begin{block}{Check}
      TODO: the answer, with the one step that decides it.
    \end{block}}
\end{frame}
% ── end of cycle ──────────────────────────────────────────────────────────────

% WRAP-UP — the definitions in one frame (what the cover's Keep in Mind says) plus the target
% misconception as a caution.
\begin{frame}[t, plain]
  \forestheader{Wrap-up: what to keep}
  \sectionlabel{5 minutes}
  \begin{itemize}\setlength{\itemsep}{5pt}
    \item TODO: definition 1, one line.
  \end{itemize}
  \vspace{4pt}
  \begin{block}{Watch out}
    TODO: the target misconception, stated as a caution.
  \end{block}
\end{frame}

% HOMEWORK — author ONE of the three blocks, per the lesson's homework source (LESSON_SHAPE.md §2).
\begin{frame}[t, plain]
  \forestheader{Homework --- start it now}
  \sectionlabel[goldacc]{10 minutes in class, alone}
  % (a) Generated (the packet page):
  \begin{block}{Homework --- scored, due the first class after two study halls}
    The \textbf{Homework} pages in your packet. Start item~1 now; the rest is due the first class
    after two study halls. TODO: what it covers.
  \end{block}
  % (b) DeltaMath, online:
  % \begin{block}{DeltaMath --- scored, due the first class after two study halls}
  %   Assignment: \textbf{TODO set name}. Start it now on your device.
  % \end{block}
  % (c) DeltaMath, printed: as (a), naming the DeltaMath pages in your packet.
  \vspace{4pt}
  \textbf{Next:} TODO one-line preview.
\end{frame}

\end{document}
